%% ============================================================
%% PARTE IV — AVALIAÇÕES DIAGNÓSTICAS
%% ============================================================
%% Código da avaliação | Nível alvo | Habilidades BNCC | Descritores CAEd
%% R-04 | Pré-silábico      | EF01LP01–EF01LP03 | D1–D4
%% R-05 | Silábico           | EF01LP02–EF01LP04 | D1–D5
%% R-06 | Silábico-Alfabético| EF01LP03–EF01LP08 | D1–D5
%% R-07 | Nível Variado      | EF01LP01–EF01LP09 | D1–D5
%% ============================================================

\chapter{Avaliações Diagnósticas}

%% ----------------------------------------------------------------
\section{Princípios da Avaliação no Método Híbrido Fonico-Kumon}
%% ----------------------------------------------------------------

A avaliação diagnóstica é componente essencial do processo de alfabetização.
No método híbrido Fonico-Kumon com Técnica da Boquinha, ela possui
características próprias que a distinguem de avaliações somativas tradicionais:

\begin{enumerate}[leftmargin=2cm]
  \item \textbf{Formativa}: acompanha o processo de aprendizagem sem punir o erro;
        cada equívoco é fonte de informação para o planejamento didático.
  \item \textbf{Diagnóstica}: identifica o nível psicogenético atual do aluno
        (pré-silábico, silábico, silábico-alfabético ou alfabético) e orienta
        a sequência de ensino.
  \item \textbf{Contínua}: é aplicada a cada 4 semanas, permitindo mapear a
        evolução ao longo do ano letivo.
  \item \textbf{Rubrica clara}: critérios bem definidos para cada nível de
        desempenho, garantindo transparência para professor, aluno e família.
  \item \textbf{Multidimensional}: avalia reconhecimento, formação de sílabas,
        leitura e produção escrita emna mesma avaliação, oferecendo um perfil
        completo do estudante.
\end{enumerate}

%% ----------------------------------------------------------------
\section{Tabela de Níveis de Alfabetização}
%% ----------------------------------------------------------------

A tabela a seguir apresenta os quatro níveis psicogenéticos de alfabetização,
conforme Ferreiro e Teberosky (1986), com a codificação BNCC e as ações
pedagógicas correspondentes. Os códigos de habilidade EF01LP referem-se
às competências de Língua Portuguesa do 1º ano do Ensino Fundamental
estabelecidas pela Base Nacional Comum Curricular (BRASIL, 2018).

\begin{table}[h]
\centering
\caption{Níveis de Alfabetização — Classificação e Ações Pedagógicas}
\label{tab:niveis-alfabetizacao}
\renewcommand{\arraystretch}{1.35}
\begin{tabular}{|c|c|p{3.2cm}|p{3.8cm}|p{3.8cm}|}
\hline
\textbf{N\'{i}vel} &
\textbf{C\'{o}d. BNCC} &
\textbf{Caracter\'{i}sticas} &
\textbf{Descritores CAEd} &
\textbf{A\c{c}\~{a}o Pedag\'{o}gica} \\
\hline
\rowcolor{warning!15}
Pré-silábico &
EF01LP01 &
O aluno ainda n\~{a}o compreende que as letras representam sons da fala.
Reconhece letras isoladamente, mas n\~{a}o estabelece rela\c{c}\~{a}o
fonema--grafema. &
D1: Identifica fonemas em s\'{i}labas \\
\hspace{0.5em} &
\hspace{0.5em} &
\hspace{0.5em} &
D2: Reconhece letras como representação de sons &
Refor\c{c}o intensivo de consci\^{e}ncia fonol\'{o}gica; atividades de
rima, alitera\c{c}\~{a}o e segmenta\c{c}\~{a}o de s\'{i}labas orais;
jogos de associa\c{c}\~{a}o imagem--letra. \\
\hline
\rowcolor{accent!10}
Silábico &
EF01LP02 &
Compreende que letras representam sons, mas ainda n\~{a}o domina a
rela\c{c}\~{a}o exata entre fonema e grafema em todas as posi\c{c}\~{o}es.
Lê s\'{i}labas isoladas com hesita\c{c}\~{a}o. &
D1: Identifica fonemas \\
\hspace{0.5em} &
\hspace{0.5em} &
\hspace{0.5em} &
D2: Forma s\'{i}labas diretas &
Consolidar s\'{i}labas diretas (CV); exerc\'{i}cios de decodifica\c{c}\~{a}o
com palavras de alta frequ\^{e}ncia; leitura de cartilhas com vocabul\'{a}rio
controlado. \\
\hline
\rowcolor{primary!10}
Silábico-Alfabético &
EF01LP03--EF01LP04 &
Utiliza estrat\'{e}gias mistas: decodifica parcialmente e utiliza pistas
do contexto. L\'{e} palavras simples com apoio. &
D1--D3: Leitura e produ\c{c}\~{a}o de s\'{i}labas e palavras simples &
Integrar leitura e produ\c{c}\~{a}o escrita; atividades de
completa\c{c}\~{a}o de frases; jogos de constru\c{c}\~{a}o de palavras
com manipul\'{a}veis. \\
\hline
\rowcolor{success!10}
Alfabético &
EF01LP05--EF01LP09 &
Domina o princ\'{i}pio alfabético: relaciona fonemas a grafemas em
todas as posi\c{c}\~{o}es. L\'{e} palavras e frases com flu\^{e}ncia. &
D1--D5: Leitura e produ\c{c}\~{a}o de frases e textos simples &
Avan\c{c}ar para leitura de textos aut\^{e}nticos; produ\c{c}\~{a}o
escrita com intencionalidade; pr\'{a}ticas de letramento. \\
\hline
\end{tabular}
\end{table}

%% ----------------------------------------------------------------
\section{Rubrica de Avaliação por Habilidade}
%% ----------------------------------------------------------------

A rubrica a seguir padroniza a avalia\c{c}\~{a}o de cada habilidade
avaliada nas quatro se\c{c}\~{o}es das avalia\c{c}\~{o}es diagnósticas.
Os n\'{i}veis s\~{a}o: \textbf{Iniciante} (0--3 pontos),
\textbf{Em desenvolvimento} (4--6 pontos) e \textbf{Dom\'{i}nio} (7--10 pontos).

\begin{table}[h]
\centering
\caption{Rubrica de Avaliação — N\'{i}veis de Desempenho por Habilidade}
\label{tab:rubrica}
\renewcommand{\arraystretch}{1.3}
\begin{tabular}{|p{3cm}|p{3.5cm}|p{3.5cm}|p{3.5cm}|}
\hline
\textbf{Habilidade} &
\textbf{N\'{i}vel 1 — Iniciante (0--3)} &
\textbf{N\'{i}vel 2 — Em desenvolvimento (4--6)} &
\textbf{N\'{i}vel 3 — Dom\'{i}nio (7--10)} \\
\hline
Reconhecimento visual de letras &
N\~{a}o identifica letras ou confunde a maioria das letras com
figuras. &
Identifica 3--4 letras corretamente; confunde letras visualmente
semelhantes (b/d, p/q). &
Identifica todas as letras do alfabeto; diferencia maiúsculas e
minúsculas. \\
\hline
Produ\c{c}\~{a}o gr\'{a}fica das letras &
N\~{a}o escreve letras ou produz tra\c{c}os
n\~{a}o identific\'{a}veis. &
Escreve 2--3 letras corretamente; tra\c{c}os com algum
controle motor. &
Escreve todas as letras com tra\c{c}ado correto e
tamanho adequado. \\
\hline
Associa\c{c}\~{a}o fonema--grafema &
N\~{a}o associa sons a letras; lê por
memoriza\c{c}\~{a}o ou adivinha\c{c}\~{a}o. &
Associa parcialmente: reconhece 2--3
correspond\^{e}ncias fonema--grafema. &
Associa todas as correspond\^{e}ncias estudadas com
precis\~{a}o. \\
\hline
Forma\c{c}\~{a}o de s\'{i}labas &
N\~{a}o forma s\'{i}labas; n\~{a}o compreende a
estrutura sil\'{a}bica. &
Forma s\'{i}labas com ajuda do professor ou
de manipul\'{a}veis. &
Forma s\'{i}labas independentemente, incluindo
compostas (ex.: TRA, PRA). \\
\hline
Separa\c{c}\~{a}o sil\'{a}bica &
N\~{a}o separa palavras em s\'{i}labas. &
Separa palavras simples com ajuda;
confunde divis\~{a}o em palavras compostas. &
Separa palavras de qualquer extens\~{a}o
independentemente. \\
\hline
Leitura de palavras &
N\~{a}o l\'{e} palavras; soletra letra por
letra sem sentido. &
L\'{e} com soletra\c{c}\~{a}o lenta (letreando);
compreende o significado ap\'{o}s decodificar. &
L\'{e} com flu\^{e}ncia; reconhece palavras de
alta frequ\^{e}ncia automaticamente. \\
\hline
Leitura de frases &
N\~{a}o l\'{e} frases; n\~{a}o compreende a
fun\c{c}\~{a}o do ponto final. &
L\'{e} frases simples com pontuação; compreende
o sentido geral. &
L\'{e} frases com entona\c{c}\~{a}o adequada;
identifica sujeito e predicado. \\
\hline
Produ\c{c}\~{a}o de frases &
N\~{a}o escreve frases; produz s\'{i}labas ou
palavras isoladas. &
Escreve frases simples com ajuda do professor
ou de modelo. &
Escreve frases independentemente, com pontua\c{c}\~{a}o
b\'{a}sica e sentido completo. \\
\hline
\end{tabular}
\end{table}

%% ============================================================
%% AVALIAÇÃO 1 — UNIDADE 1: VOGAIS
%% Código: R-04 | Nível: Pré-silábico | BNCC: EF01LP01–EF01LP03
%% Descritores: D1 (identificação de fonemas), D2 (reconhecimento de
%% letras como representação de sons), D3 (associação fonema–grafema),
%% D4 (produção gráfica)
%% ============================================================

\section{Avaliação da Unidade 1 --- Vogais}

\begin{avaliacao}
%% --- BarrA DE METADADOS ---
\begin{center}
\resizebox{\textwidth}{!}{%
\begin{tabular}{@{}llll@{}}
\textbf{Código:} R-04 &
\textbf{N\'{i}vel:} Pré-silábico &
\textbf{BNCC:} EF01LP01--EF01LP03 &
\textbf{Descritores:} D1--D4 \\[0.3em]
\textbf{Unidade:} 1 — Vogais &
\textbf{Pontuação:} 40 pontos &
\textbf{Tempo:} 40 min &
\textbf{Frequência:} A cada 4 semanas \\
\end{tabular}}
\end{center}

\vspace{0.5em}
\noindent\textbf{Nome:} \underline{\hspace{5cm}} \quad
\textbf{Escola:} \underline{\hspace{4cm}} \quad
\textbf{Data:} \underline{\hspace{2.5cm}}\\[0.5em]
\noindent\textbf{Turma:} \underline{\hspace{2cm}} \quad
\textbf{N\'{i}vel Atual:} \underline{\hspace{3cm}}

%% --- SEÇÃO 1: RECONHECIMENTO (10 pontos) ---
\subsection*{Seção 1 — Reconhecimento de Letras (10 pontos)}

\textit{Objetivo:} avaliar se o aluno reconhece visualmente as vogais
e compreende que cada letra corresponde a um som da fala.

\begin{enumerate}
  \item \textbf{Escrever cada vogal} (1 ponto cada — total 5 pontos):

  \begin{center}
  \begin{tabular}{ccccc}
    A: \underline{\hspace{1.5cm}} &
    E: \underline{\hspace{1.5cm}} &
    I: \underline{\hspace{1.5cm}} &
    O: \underline{\hspace{1.5cm}} &
    U: \underline{\hspace{1.5cm}} \\
  \end{tabular}
  \end{center}

  \item \textbf{Identificar a vogal que falta na palavra} (1 ponto cada — total 5 pontos):

  \begin{center}
  \begin{tabular}{cccc}
    \underline{\hspace{0.8cm}}sa &
    m\underline{\hspace{0.8cm}}sa &
    s\underline{\hspace{0.8cm}}po &
    l\underline{\hspace{0.8cm}}va \\[0.3em]
    ( \textbf{A} ) &
    ( \textbf{E} ) &
    ( \textbf{I} ) &
    ( \textbf{U} ) \\
  \end{tabular}
  \end{center}

  \item \textbf{Circundar as letras que o professor fala} (total incluído nos 10 pontos):
  o professor pronuncia 3 vogais; o aluno deve circundar as letras correspondentes
  em uma folha com todas as 5 vogais impressas.
\end{enumerate}

\medskip
\noindent\textbf{Nota da Seção 1:} \underline{\hspace{1cm}} / 10

%% --- SEÇÃO 2: FORMAÇÃO DE SÍLABAS (10 pontos) ---
\subsection*{Seção 2 — Formação de Sílabas (10 pontos)}

\textit{Objetivo:} verificar se o aluno compreende a estrutura silábica
e é capaz de combinar consoante + vogal para formar sílabas diretas.

\begin{enumerate}[resume]
  \item \textbf{Formar sílabas unindo consoante e vogal} (1 ponto cada — total 5 pontos):

  \begin{center}
  \resizebox{\textwidth}{!}{%
\begin{tabular}{ccccc}
    \sil{M}{A} & \sil{S}{E} & \sil{T}{I} & \sil{R}{O} & \sil{L}{U} \\[0.5em]
    \sil{B}{A} & \sil{C}{E} & \sil{D}{I} & \sil{F}{O} & \sil{G}{U} \\
  \end{tabular}}
  \end{center}

  \item \textbf{Separar palavras em sílabas} (1 ponto cada — total 5 pontos):

  \begin{center}
  \resizebox{\textwidth}{!}{%
\begin{tabular}{ccc}
    MESA $\rightarrow$ \underline{\hspace{1.2cm}} + \underline{\hspace{1.2cm}} &
    CASA $\rightarrow$ \underline{\hspace{1.2cm}} + \underline{\hspace{1.2cm}} &
    SAPO $\rightarrow$ \underline{\hspace{1.2cm}} + \underline{\hspace{1.2cm}} \\[0.5em]
    ( ME + SA ) &
    ( CA + SA ) &
    ( SA + PO ) \\
  \end{tabular}}
  \end{center}
\end{enumerate}

\medskip
\noindent\textbf{Nota da Seção 2:} \underline{\hspace{1cm}} / 10

%% --- SEÇÃO 3: LEITURA (10 pontos) ---
\subsection*{Seção 3 — Leitura (10 pontos)}

\textit{Objetivo:} avaliar se o aluno consegue decodificar palavras e frases
simples, aplicando o conhecimento das vogais e das sílabas diretas.

\begin{enumerate}[resume]
  \item \textbf{Ler palavras} (1 ponto cada — total 5 pontos):

  \begin{center}
  $\rightarrow$ mesa \quad casa \quad sapo \quad luto \quad resto
  \end{center}

  \item \textbf{Ler frases} (1 ponto cada — total 5 pontos):

  \begin{enumerate}[label=\arabic*., leftmargin=1.5cm]
    \item O gato mia.
    \item A menina pula.
    \item O menino corre.
  \end{enumerate}
\end{enumerate}

\medskip
\noindent\textbf{Nota da Seção 3:} \underline{\hspace{1cm}} / 10

%% --- SEÇÃO 4: PRODUÇÃO (10 pontos) ---
\subsection*{Seção 4 — Produção Escrita (10 pontos)}

\textit{Objetivo:} verificar se o aluno é capaz de produzir escrita
espontânea com intencionalidade, usando as vogais e sílabas estudadas.

\begin{enumerate}[resume]
  \item \textbf{Escrever 3 palavras com cada vogal} (1 ponto por vogal — total 5 pontos):

  \begin{center}
  \begin{tabular}{l}
    A: \underline{\hspace{2.5cm}} \quad \underline{\hspace{2.5cm}} \quad \underline{\hspace{2.5cm}} \\[0.5em]
    E: \underline{\hspace{2.5cm}} \quad \underline{\hspace{2.5cm}} \quad \underline{\hspace{2.5cm}} \\[0.5em]
    I: \underline{\hspace{2.5cm}} \quad \underline{\hspace{2.5cm}} \quad \underline{\hspace{2.5cm}} \\[0.5em]
    O: \underline{\hspace{2.5cm}} \quad \underline{\hspace{2.5cm}} \quad \underline{\hspace{2.5cm}} \\[0.5em]
    U: \underline{\hspace{2.5cm}} \quad \underline{\hspace{2.5cm}} \quad \underline{\hspace{2.5cm}} \\
  \end{tabular}
  \end{center}

  \item \textbf{Escrever uma frase simples} (5 pontos):
  \begin{center}
  \underline{\hspace{10cm}}
  \end{center}
  \textit{(Nota: 5 pontos = frase completa com sujeito + verbo;
  3 pontos = frase incompleta ou com erros graves;
  1 ponto = palavras isoladas sem relação.)}
\end{enumerate}

\medskip
\noindent\textbf{Nota da Seção 4:} \underline{\hspace{1cm}} / 10

%% --- NOTA TOTAL ---
\begin{center}
\large\textbf{NOTA TOTAL:} \underline{\hspace{1.5cm}} / 40
\end{center}

\begin{center}
\textbf{Classificação:} \underline{\hspace{3cm}}
\hspace{1cm}
\textbf{Nível de Alfabetização:} \underline{\hspace{3cm}}
\end{center}

\end{avaliacao}

%% ============================================================
%% AVALIAÇÃO 2 — UNIDADE 2: CONSONANTES
%% Código: R-05 | Nível: Silábico | BNCC: EF01LP02–EF01LP04
%% Descritores: D1 (identificação de fonemas), D2 (formação de sílabas),
%% D3 (leitura de palavras), D4 (separação silábica),
%% D5 (produção de frases)
%% ============================================================

\section{Avaliação da Unidade 2 --- Consonantes do Bloco 1}

\begin{avaliacao}
%% --- BARRA DE METADADOS ---
\begin{center}
\resizebox{\textwidth}{!}{%
\begin{tabular}{@{}llll@{}}
\textbf{Código:} R-05 &
\textbf{N\'{i}vel:} Sil\'{a}bico &
\textbf{BNCC:} EF01LP02--EF01LP04 &
\textbf{Descritores:} D1--D5 \\[0.3em]
\textbf{Unidade:} 2 — Consonantes (M, S, T, R, L) &
\textbf{Pontua\c{c}\~{a}o:} 40 pontos &
\textbf{Tempo:} 40 min &
\textbf{Frequ\^{e}ncia:} A cada 4 semanas \\
\end{tabular}}
\end{center}

\vspace{0.5em}
\noindent\textbf{Nome:} \underline{\hspace{5cm}} \quad
\textbf{Escola:} \underline{\hspace{4cm}} \quad
\textbf{Data:} \underline{\hspace{2.5cm}}\\[0.5em]
\noindent\textbf{Turma:} \underline{\hspace{2cm}} \quad
\textbf{N\'{i}vel Atual:} \underline{\hspace{3cm}}

%% --- SEÇÃO 1: RECONHECIMENTO (10 pontos) ---
\subsection*{Seção 1 — Reconhecimento de Consonantes (10 pontos)}

\textit{Objetivo:} avaliar se o aluno reconhece as consonantes M, S, T, R e L
e identifica sua posição nas palavras.

\begin{enumerate}
  \item \textbf{Identificar a consonante que inicia cada palavra} (1 ponto cada — total 5 pontos):

  \begin{center}
  \begin{tabular}{ccccc}
    \underline{\hspace{0.8cm}}esa &
    \underline{\hspace{0.8cm}}apo &
    \underline{\hspace{0.8cm}}eto &
    \underline{\hspace{0.8cm}}ato &
    \underline{\hspace{0.8cm}}ua \\[0.3em]
    ( \textbf{M} ) &
    ( \textbf{S} ) &
    ( \textbf{T} ) &
    ( \textbf{R} ) &
    ( \textbf{L} ) \\
  \end{tabular}
  \end{center}

  \item \textbf{Circular as palavras que começam com ``M''} (1 ponto por acerto — total 5 pontos):

  \begin{center}
  \begin{tabular}{ccccc}
    \fbox{macaco} & gato & \fbox{mesa} & sapo & \fbox{milho} \\[0.3em]
    rato & lata & \fbox{mola} & teto & \fbox{mula}
  \end{tabular}
  \end{center}
\end{enumerate}

\medskip
\noindent\textbf{Nota da Seção 1:} \underline{\hspace{1cm}} / 10

%% --- SEÇÃO 2: FORMAÇÃO DE SÍLABAS (10 pontos) ---
\subsection*{Seção 2 — Formação de Sílabas (10 pontos)}

\textit{Objetivo:} verificar se o aluno é capaz de combinar cada consonante
com as 5 vogais para formar sílabas diretas.

\begin{enumerate}[resume]
  \item \textbf{Formar sílabas com M} (1 ponto cada — total 3 pontos):

  \begin{center}
  \sil{M}{A} \quad \sil{M}{E} \quad \sil{M}{I} \quad \sil{M}{O} \quad \sil{M}{U}
  \end{center}

  \item \textbf{Formar sílabas com S} (1 ponto cada — total 3 pontos):

  \begin{center}
  \sil{S}{A} \quad \sil{S}{E} \quad \sil{S}{I} \quad \sil{S}{O} \quad \sil{S}{U}
  \end{center}

  \item \textbf{Formar sílabas com T} (1 ponto cada — total 2 pontos):

  \begin{center}
  \sil{T}{A} \quad \sil{T}{E} \quad \sil{T}{I} \quad \sil{T}{O} \quad \sil{T}{U}
  \end{center}

  \item \textbf{Formar sílabas com R} (1 ponto — total 1 ponto):

  \begin{center}
  \sil{R}{A} \quad \sil{R}{E} \quad \sil{R}{I} \quad \sil{R}{O} \quad \sil{R}{U}
  \end{center}

  \item \textbf{Formar sílabas com L} (1 ponto — total 1 ponto):

  \begin{center}
  \sil{L}{A} \quad \sil{L}{E} \quad \sil{L}{I} \quad \sil{L}{O} \quad \sil{L}{U}
  \end{center}
\end{enumerate}

\medskip
\noindent\textbf{Nota da Seção 2:} \underline{\hspace{1cm}} / 10

%% --- SEÇÃO 3: LEITURA (10 pontos) ---
\subsection*{Seção 3 — Leitura (10 pontos)}

\textit{Objetivo:} avaliar a capacidade de decodificação de palavras
com as consonantes estudadas, em contexto de frases simples.

\begin{enumerate}[resume]
  \item \textbf{Ler palavras com M} (1 ponto cada — total 3 pontos):

  \begin{center}
  $\rightarrow$ mesa \quad macaco \quad milho \quad mola \quad mula
  \end{center}

  \item \textbf{Ler palavras com S} (1 ponto cada — total 2 pontos):

  \begin{center}
  $\rightarrow$ sapo \quad sol \quad sapato \quad sola \quad suco
  \end{center}

  \item \textbf{Ler frases simples} (1 ponto cada — total 5 pontos):

  \begin{enumerate}[label=\arabic*., leftmargin=1.5cm]
    \item O sapo salta.
    \item A mesa é redonda.
    \item O menino lê um livro.
  \end{enumerate}
\end{enumerate}

\medskip
\noindent\textbf{Nota da Seção 3:} \underline{\hspace{1cm}} / 10

%% --- SEÇÃO 4: PRODUÇÃO (10 pontos) ---
\subsection*{Seção 4 — Produção Escrita (10 pontos)}

\textit{Objetivo:} verificar se o aluno produz palavras e frases
usando as consonantes e sílabas diretas estudadas.

\begin{enumerate}[resume]
  \item \textbf{Escrever 2 palavras com cada consonante} (1 ponto cada — total 5 pontos):

  \begin{center}
  \begin{tabular}{l}
    M: \underline{\hspace{2.5cm}} \quad \underline{\hspace{2.5cm}} \\[0.5em]
    S: \underline{\hspace{2.5cm}} \quad \underline{\hspace{2.5cm}} \\[0.5em]
    T: \underline{\hspace{2.5cm}} \quad \underline{\hspace{2.5cm}} \\[0.5em]
    R: \underline{\hspace{2.5cm}} \quad \underline{\hspace{2.5cm}} \\[0.5em]
    L: \underline{\hspace{2.5cm}} \quad \underline{\hspace{2.5cm}} \\
  \end{tabular}
  \end{center}

  \item \textbf{Escrever 2 frases usando palavras com as consonantes estudadas}
  (2,5 pontos cada — total 5 pontos):

  \begin{center}
  \underline{\hspace{10cm}} \\[0.8em]
  \underline{\hspace{10cm}}
  \end{center}
\end{enumerate}

\medskip
\noindent\textbf{Nota da Seção 4:} \underline{\hspace{1cm}} / 10

%% --- NOTA TOTAL ---
\begin{center}
\large\textbf{NOTA TOTAL:} \underline{\hspace{1.5cm}} / 40
\end{center}

\begin{center}
\textbf{Classificação:} \underline{\hspace{3cm}}
\hspace{1cm}
\textbf{N\'{i}vel de Alfabetiza\c{c}\~{a}o:} \underline{\hspace{3cm}}
\end{center}

\end{avaliacao}

%% ============================================================
%% AVALIAÇÃO 3 — UNIDADE 3: SÍLABAS
%% Código: R-06 | Nível: Silábico-Alfabético | BNCC: EF01LP03–EF01LP08
%% Descritores: D1 (identificação de fonemas), D2 (formação de sílabas),
%% D3 (leitura de palavras), D4 (separação silábica),
%% D5 (produção de frases)
%% ============================================================

\section{Avaliação da Unidade 3 --- Sílabas Diretas}

\begin{avaliacao}
%% --- BARRA DE METADADOS ---
\begin{center}
\resizebox{\textwidth}{!}{%
\begin{tabular}{@{}llll@{}}
\textbf{Código:} R-06 &
\textbf{N\'{i}vel:} Sil\'{a}bico-Alfab\'{e}tico &
\textbf{BNCC:} EF01LP03--EF01LP08 &
\textbf{Descritores:} D1--D5 \\[0.3em]
\textbf{Unidade:} 3 — S\'{i}labas Diretas &
\textbf{Pontua\c{c}\~{a}o:} 40 pontos &
\textbf{Tempo:} 40 min &
\textbf{Frequ\^{e}ncia:} A cada 4 semanas \\
\end{tabular}}
\end{center}

\vspace{0.5em}
\noindent\textbf{Nome:} \underline{\hspace{5cm}} \quad
\textbf{Escola:} \underline{\hspace{4cm}} \quad
\textbf{Data:} \underline{\hspace{2.5cm}}\\[0.5em]
\noindent\textbf{Turma:} \underline{\hspace{2cm}} \quad
\textbf{N\'{i}vel Atual:} \underline{\hspace{3cm}}

%% --- SEÇÃO 1: RECONHECIMENTO (10 pontos) ---
\subsection*{Seção 1 — Reconhecimento e Formação de Sílabas (10 pontos)}

\textit{Objetivo:} avaliar se o aluno domina a forma\c{c}\~{a}o de
s\'{i}labas diretas com todas as consonantes e vogais estudadas.

\begin{enumerate}
  \item \textbf{Formar s\'{i}labas com todas as vogais} (1 ponto cada — total 10 pontos):

  \begin{center}
  \resizebox{\textwidth}{!}{%
\begin{tabular}{ccccc}
    \sil{M}{A} & \sil{M}{E} & \sil{M}{I} & \sil{M}{O} & \sil{M}{U} \\
    \sil{S}{A} & \sil{S}{E} & \sil{S}{I} & \sil{S}{O} & \sil{S}{U} \\
    \sil{T}{A} & \sil{T}{E} & \sil{T}{I} & \sil{T}{O} & \sil{T}{U} \\
    \sil{R}{A} & \sil{R}{E} & \sil{R}{I} & \sil{R}{O} & \sil{R}{U} \\
    \sil{L}{A} & \sil{L}{E} & \sil{L}{I} & \sil{L}{O} & \sil{L}{U} \\
  \end{tabular}}
  \end{center}
\end{enumerate}

\medskip
\noindent\textbf{Nota da Seção 1:} \underline{\hspace{1cm}} / 10

%% --- SEÇÃO 2: SEPARAÇÃO SILÁBICA (10 pontos) ---
\subsection*{Seção 2 — Separação Silábica (10 pontos)}

\textit{Objetivo:} verificar se o aluno consegue separar palavras em
sílabas, identificando a fronteira silábica corretamente.

\begin{enumerate}[resume]
  \item \textbf{Separar palavras em sílabas} (1 ponto cada — total 10 pontos):

  \begin{center}
  \resizebox{\textwidth}{!}{%
\begin{tabular}{ccc}
    MESA $\rightarrow$ \underline{\hspace{1.2cm}} + \underline{\hspace{1.2cm}} &
    CASA $\rightarrow$ \underline{\hspace{1.2cm}} + \underline{\hspace{1.2cm}} &
    SAPO $\rightarrow$ \underline{\hspace{1.2cm}} + \underline{\hspace{1.2cm}} \\[0.8em]
    RESTO $\rightarrow$ \underline{\hspace{1.2cm}} + \underline{\hspace{1.2cm}} &
    MOLA $\rightarrow$ \underline{\hspace{1.2cm}} + \underline{\hspace{1.2cm}} &
    TIRSO $\rightarrow$ \underline{\hspace{1.2cm}} + \underline{\hspace{1.2cm}} \\[0.8em]
    SALA $\rightarrow$ \underline{\hspace{1.2cm}} + \underline{\hspace{1.2cm}} &
    LUTO $\rightarrow$ \underline{\hspace{1.2cm}} + \underline{\hspace{1.2cm}} &
    RATO $\rightarrow$ \underline{\hspace{1.2cm}} + \underline{\hspace{1.2cm}} \\[0.8em]
    MESA $\rightarrow$ \underline{\hspace{1.2cm}} + \underline{\hspace{1.2cm}} &
    TIRA $\rightarrow$ \underline{\hspace{1.2cm}} + \underline{\hspace{1.2cm}} &
    ROLA $\rightarrow$ \underline{\hspace{1.2cm}} + \underline{\hspace{1.2cm}} \\
  \end{tabular}}
  \end{center}
\end{enumerate}

\medskip
\noindent\textbf{Nota da Seção 2:} \underline{\hspace{1cm}} / 10

%% --- SEÇÃO 3: LEITURA (10 pontos) ---
\subsection*{Seção 3 — Leitura (10 pontos)}

\textit{Objetivo:} avaliar a fluência de leitura de palavras e frases,
considerando a velocidade e a compreensão do sentido.

\begin{enumerate}[resume]
  \item \textbf{Ler palavras} (1 ponto cada — total 5 pontos):

  \begin{center}
  $\rightarrow$ mesa \quad sala \quad luto \quad resto \quad mola
  \end{center}

  \item \textbf{Ler palavras mais desafiadoras} (1 ponto cada — total 3 pontos):

  \begin{center}
  $\rightarrow$ tira \quad rola \quad dica \quad mula \quad suco
  \end{center}

  \item \textbf{Ler frases} (1 ponto cada — total 2 pontos):

  \begin{enumerate}[label=\arabic*., leftmargin=1.5cm]
    \item O rato corre na sala.
    \item A mola é vermelha.
  \end{enumerate}
\end{enumerate}

\medskip
\noindent\textbf{Nota da Seção 3:} \underline{\hspace{1cm}} / 10

%% --- SEÇÃO 4: PRODUÇÃO (10 pontos) ---
\subsection*{Seção 4 — Produção Escrita (10 pontos)}

\textit{Objetivo:} avaliar se o aluno produz escrita com intencionalidade,
usando palavras e frases com as sílabas diretas estudadas.

\begin{enumerate}[resume]
  \item \textbf{Escrever 4 palavras que o professor ditado} (2 pontos cada — total 4 pontos):

  \begin{center}
  (O professor dita: \textbf{mesa, sapo, rato, luto})
  \end{center}

  \item \textbf{Separar as palavras ditadas em sílabas} (1 ponto cada — total 3 pontos):

  \begin{center}
  \resizebox{\textwidth}{!}{%
\begin{tabular}{ccc}
    \underline{\hspace{1.5cm}} + \underline{\hspace{1.5cm}} &
    \underline{\hspace{1.5cm}} + \underline{\hspace{1.5cm}} &
    \underline{\hspace{1.5cm}} + \underline{\hspace{1.5cm}} \\[0.5em]
    ( mesa ) & ( sapo ) & ( rato ) \\
  \end{tabular}}
  \end{center}

  \item \textbf{Escrever 2 frases usando palavras com sílabas diretas}
  (1,5 pontos cada — total 3 pontos):

  \begin{center}
  \underline{\hspace{10cm}} \\[0.8em]
  \underline{\hspace{10cm}}
  \end{center}
\end{enumerate}

\medskip
\noindent\textbf{Nota da Seção 4:} \underline{\hspace{1cm}} / 10

%% --- NOTA TOTAL ---
\begin{center}
\large\textbf{NOTA TOTAL:} \underline{\hspace{1.5cm}} / 40
\end{center}

\begin{center}
\textbf{Classificação:} \underline{\hspace{3cm}}
\hspace{1cm}
\textbf{N\'{i}vel de Alfabetiza\c{c}\~{a}o:} \underline{\hspace{3cm}}
\end{center}

\end{avaliacao}

%% ============================================================
%% AVALIAÇÃO FINAL — 1º ANO
%% Código: R-07 | Nível: Variado | BNCC: EF01LP01–EF01LP09
%% Descritores: D1 (identificação de fonemas), D2 (formação de sílabas),
%% D3 (leitura de palavras), D4 (separação silábica),
%% D5 (produção de frases)
%% ============================================================

\section{Avaliação Final do 1º Ano}

\begin{avaliacao}
%% --- BARRA DE METADADOS ---
\begin{center}
\resizebox{\textwidth}{!}{%
\begin{tabular}{@{}llll@{}}
\textbf{C\'{o}digo:} R-07 &
\textbf{N\'{i}vel:} Variado &
\textbf{BNCC:} EF01LP01--EF01LP09 &
\textbf{Descritores:} D1--D5 \\[0.3em]
\textbf{Avalia\c{c}\~{a}o:} Final do 1º Ano &
\textbf{Pontua\c{c}\~{a}o:} 40 pontos &
\textbf{Tempo:} 50 min &
\textbf{Frequ\^{e}ncia:} Semestral \\
\end{tabular}}
\end{center}

\vspace{0.5em}
\noindent\textbf{Nome:} \underline{\hspace{5cm}} \quad
\textbf{Escola:} \underline{\hspace{4cm}} \quad
\textbf{Data:} \underline{\hspace{2.5cm}}\\[0.5em]
\noindent\textbf{Turma:} \underline{\hspace{2cm}} \quad
\textbf{N\'{i}vel Atual:} \underline{\hspace{3cm}}

%% --- SEÇÃO 1: RECONHECIMENTO (10 pontos) ---
\subsection*{Seção 1 — Reconhecimento de Letras e Sílabas (10 pontos)}

\textit{Objetivo:} avaliar o reconhecimento global do alfabeto e a
compreensão das relações fonema--grafema em sílabas diretas e compostas.

\begin{enumerate}
  \item \textbf{Escrever cada vogal} (1 ponto cada — total 5 pontos):

  \begin{center}
  \begin{tabular}{ccccc}
    A: \underline{\hspace{1.5cm}} &
    E: \underline{\hspace{1.5cm}} &
    I: \underline{\hspace{1.5cm}} &
    O: \underline{\hspace{1.5cm}} &
    U: \underline{\hspace{1.5cm}} \\
  \end{tabular}
  \end{center}

  \item \textbf{Escrever cada consonante do Bloco 1} (1 ponto cada — total 5 pontos):

  \begin{center}
  \begin{tabular}{ccccc}
    M: \underline{\hspace{1.5cm}} &
    S: \underline{\hspace{1.5cm}} &
    T: \underline{\hspace{1.5cm}} &
    R: \underline{\hspace{1.5cm}} &
    L: \underline{\hspace{1.5cm}} \\
  \end{tabular}
  \end{center}
\end{enumerate}

\medskip
\noindent\textbf{Nota da Seção 1:} \underline{\hspace{1cm}} / 10

%% --- SEÇÃO 2: FORMAÇÃO DE SÍLABAS (10 pontos) ---
\subsection*{Seção 2 — Formação de Sílabas (10 pontos)}

\textit{Objetivo:} verificar se o aluno domina a forma\c{c}\~{a}o de
s\'{i}labas diretas e compreende a estrutura sil\'{a}bica de palavras.

\begin{enumerate}[resume]
  \item \textbf{Formar s\'{i}labas} (1 ponto cada — total 5 pontos):

  \begin{center}
  \resizebox{\textwidth}{!}{%
\begin{tabular}{ccccc}
    \sil{M}{A} & \sil{S}{E} & \sil{T}{I} & \sil{R}{O} & \sil{L}{U} \\
    \sil{B}{A} & \sil{C}{E} & \sil{D}{I} & \sil{F}{O} & \sil{G}{U} \\
  \end{tabular}}
  \end{center}

  \item \textbf{Separar palavras em sílabas} (1 ponto cada — total 5 pontos):

  \begin{center}
  \resizebox{\textwidth}{!}{%
\begin{tabular}{ccc}
    MESA $\rightarrow$ \underline{\hspace{1.2cm}} + \underline{\hspace{1.2cm}} &
    CASA $\rightarrow$ \underline{\hspace{1.2cm}} + \underline{\hspace{1.2cm}} &
    SAPO $\rightarrow$ \underline{\hspace{1.2cm}} + \underline{\hspace{1.2cm}} \\[0.8em]
    RESTO $\rightarrow$ \underline{\hspace{1.2cm}} + \underline{\hspace{1.2cm}} &
    MOLA $\rightarrow$ \underline{\hspace{1.2cm}} + \underline{\hspace{1.2cm}} &
    TIRSO $\rightarrow$ \underline{\hspace{1.2cm}} + \underline{\hspace{1.2cm}} \\
  \end{tabular}}
  \end{center}
\end{enumerate}

\medskip
\noindent\textbf{Nota da Seção 2:} \underline{\hspace{1cm}} / 10

%% --- SEÇÃO 3: LEITURA (10 pontos) ---
\subsection*{Seção 3 — Leitura (10 pontos)}

\textit{Objetivo:} avaliar a fluência de leitura em três níveis de
complexidade: palavras isoladas, palavras em frase e frases completas.

\begin{enumerate}[resume]
  \item \textbf{Ler palavras simples} (1 ponto cada — total 4 pontos):

  \begin{center}
  $\rightarrow$ mesa \quad casa \quad sapo \quad luto
  \end{center}

  \item \textbf{Ler palavras mais complexas} (1 ponto cada — total 3 pontos):

  \begin{center}
  $\rightarrow$ resto \quad bola \quad tira \quad rola \quad dica
  \end{center}

  \item \textbf{Ler frases} (1 ponto cada — total 3 pontos):

  \begin{enumerate}[label=\arabic*., leftmargin=1.5cm]
    \item O gato mia.
    \item A menina pula.
    \item O menino corre no parque.
  \end{enumerate}
\end{enumerate}

\medskip
\noindent\textbf{Nota da Seção 3:} \underline{\hspace{1cm}} / 10

%% --- SEÇÃO 4: PRODUÇÃO (10 pontos) ---
\subsection*{Seção 4 — Produção Escrita (10 pontos)}

\textit{Objetivo:} avaliar a capacidade de produzir escrita espontânea
com intencionalidade, usando vocabulário e estruturas gramaticais adequadas.

\begin{enumerate}[resume]
  \item \textbf{Escrever 3 palavras com cada vogal} (1 ponto por vogal — total 5 pontos):

  \begin{center}
  \begin{tabular}{l}
    A: \underline{\hspace{2.5cm}} \quad \underline{\hspace{2.5cm}} \quad \underline{\hspace{2.5cm}} \\[0.5em]
    E: \underline{\hspace{2.5cm}} \quad \underline{\hspace{2.5cm}} \quad \underline{\hspace{2.5cm}} \\[0.5em]
    I: \underline{\hspace{2.5cm}} \quad \underline{\hspace{2.5cm}} \quad \underline{\hspace{2.5cm}} \\[0.5em]
    O: \underline{\hspace{2.5cm}} \quad \underline{\hspace{2.5cm}} \quad \underline{\hspace{2.5cm}} \\[0.5em]
    U: \underline{\hspace{2.5cm}} \quad \underline{\hspace{2.5cm}} \quad \underline{\hspace{2.5cm}} \\
  \end{tabular}
  \end{center}

  \item \textbf{Escrever uma frase simples} (5 pontos):

  \begin{center}
  \underline{\hspace{10cm}}
  \end{center}
  \textit{(Nota: 5 pontos = frase completa com sentido; 3 pontos = frase
  com erros de pontuação; 1 ponto = palavras isoladas.)}
\end{enumerate}

\medskip
\noindent\textbf{Nota da Seção 4:} \underline{\hspace{1cm}} / 10

%% --- NOTA TOTAL ---
\begin{center}
\large\textbf{NOTA TOTAL:} \underline{\hspace{1.5cm}} / 40
\end{center}

\begin{center}
\textbf{Classifica\c{c}\~{a}o:} \underline{\hspace{3cm}}
\hspace{1cm}
\textbf{N\'{i}vel de Alfabetiza\c{c}\~{a}o:} \underline{\hspace{3cm}}
\end{center}

\end{avaliacao}

%% ============================================================
%% TABELA CONSOLIDADA — CLASSIFICAÇÃO FINAL
%% ============================================================

\section{Classificação Final por Nível de Alfabetização}

Após a aplica\c{c}\~{a}o de todas as avalia\c{c}\~{o}es diagnósticas,
o professor deve classificar cada aluno na tabela abaixo, considerando
a pontua\c{c}\~{a}o m\'{e}dia e a evolu\c{c}\~{a}o ao longo do ano.

\begin{table}[h]
\centering
\caption{Classifica\c{c}\~{a}o do N\'{i}vel de Alfabetiza\c{c}\~{a}o --- Consolida\c{c}\~{a}o}
\label{tab:classificacao-final}
\renewcommand{\arraystretch}{1.35}
\resizebox{\textwidth}{!}{%
\begin{tabular}{|c|c|p{3cm}|p{3.8cm}|p{3.8cm}|}
\hline
\textbf{Pontua\c{c}\~{a}o} &
\textbf{N\'{i}vel} &
\textbf{C\'{o}d. BNCC} &
\textbf{Descri\c{c}\~{a}o} &
\textbf{A\c{c}\~{a}o Pedag\'{o}gica} \\
\hline
\rowcolor{warning!15}
0--10 &
Pré-sil\'{a}bico &
EF01LP01 &
N\~{a}o reconhece letras como representa\c{c}\~{a}o de sons;
confunde letras com figuras. &
Refor\c{c}o intensivo de consci\^{e}ncia fonol\'{o}gica;
atividades de rima e alitera\c{c}\~{a}o; uso de manipul\'{a}veis. \\
\hline
\rowcolor{accent!10}
11--20 &
Sil\'{a}bico &
EF01LP02 &
Reconhece que letras representam sons, mas n\~{a}o domina a
rela\c{c}\~{a}o fonema--grafema em todas as posi\c{c}\~{o}es. &
Consolidar s\'{i}labas diretas (CV); exerc\'{i}cios de
decodifica\c{c}\~{a}o com palavras de alta frequ\^{e}ncia. \\
\hline
\rowcolor{primary!10}
21--30 &
Sil\'{a}bico-Alfab\'{e}tico &
EF01LP03--EF01LP04 &
Mistura estrat\'{e}gias sil\'{a}bicas e alfab\'{e}ticas;
l\'{e} palavras simples com apoio. &
Integrar leitura e produ\c{c}\~{a}o escrita; atividades de
completa\c{c}\~{a}o de frases; jogos de constru\c{c}\~{a}o de palavras. \\
\hline
\rowcolor{success!10}
31--40 &
Alfab\'{e}tico &
EF01LP05--EF01LP09 &
Domina o princ\'{i}pio alfab\'{e}tico; l\'{e} palavras e frases
com flu\^{e}ncia. &
Avan\c{c}ar para leitura de textos aut\^{e}nticos; produ\c{c}\~{a}o
escrita com intencionalidade; pr\'{a}ticas de letramento. \\
\hline
\end{tabular}}
\end{table}

%% ============================================================
%% ORIENTAÇÕES PARA O PROFESSOR
%% ============================================================

\section{Orientações para a Aplicação das Avaliações}

\begin{enumerate}[leftmargin=2cm]
  \item \textbf{Ambiente acolhedor}: aplique a avalia\c{c}\~{a}o em
  ambiente tranquilo, sem pressa. Explique ao aluno que se trata de
  uma atividade para ajudar o professor a ensinar melhor.

  \item \textbf{Registro cuidadoso}: anote n\~{a}o apenas as notas,
  mas tamb\'{e}m as estrat\'{e}gias que o aluno utiliza (soletra\c{c}\~{a}o,
  adivinha\c{c}\~{a}o, leitura visual, etc.).

  \item \textbf{Análise do perfil}: ap\'{o}s a avalia\c{c}\~{a}o, preencha
  o quadro de classifica\c{c}\~{a}o e planeje as pr\'{o}ximas aulas
  com base no perfil do grupo e de cada aluno.

  \item \textbf{Comunica\c{c}\~{a}o com a fam\'{i}lia}: envie os
  resultados para os pais/respons\'{a}veis, explicando o n\'{i}vel
  do aluno e as a\c{c}\~{o}es pedag\'{o}gicas que ser\~{a}o adotadas.

  \item \textbf{Adapta\c{c}\~{a}o}: para alunos com necessidades
  especiais, adapte o formato da avalia\c{c}\~{a}o conforme orienta\c{c}\~{a}o
  do PAEE ou do CRP.
\end{enumerate}

%% ============================================================
%% NOVO (R201): ROTEAMENTO DIAGNÓSTICO PARA ROTAS 1A-1D
%% ============================================================
\section{Roteamento Diagnóstico para Rotas Individualizadas (1A--1D)}

O resultado da avaliação diagnóstica orienta a ativação de uma das quatro rotas
individualizadas deste livro. O critério de roteamento é o percentual de
acertos com qualidade demonstrada, nunca a velocidade de execução.

\subsection{Como Calcular o Percentual de Acertos}

\begin{enumerate}[leftmargin=2cm]
  \item Some a pontuação obtida pelo aluno em todas as seções da avaliação.
  \item Divida pela pontuação total (40 pontos) e multiplique por 100.
  \item Registre o percentual no quadro de acompanhamento da Parte 5.
\end{enumerate}

\subsection{Tabela de Roteamento}

\begin{center}
\begin{tabular}{|c|c|l|l|}
\hline
\textbf{Percentual} & \textbf{Rota} & \textbf{Ponto de Entrada} & \textbf{Nível K} \\
\hline
$<$ 50\% & 1A & Unidade 0 + folhas A-01 a A-05 & K0--K1 \\
\hline
50\%--70\% & 1B & Unidades 4--7 + folhas B-01 a C-03 & K2--K3 \\
\hline
70\%--85\% & 1C & Unidades 8--9 + folhas C-04 a D-02 & K3--K4 \\
\hline
$>$ 85\% & 1D & Unidade 10 + folhas D-03 a D-06 & K4--K5 \\
\hline
\end{tabular}
\end{center}

\subsection{Reavaliação e Migração de Rota}

A rota não é fixa. A cada reavaliação mensal ou bimestral, o professor verifica
se o aluno progrediu e, se necessário, migra de rota:

\begin{itemize}
  \item \textbf{Migração ascendente}: aluno com 90\%+ de domínio em duas
        avaliações consecutivas avança para a rota seguinte.
  \item \textbf{Permanência}: aluno com progresso parcial permanece na mesma
        rota com ajuste de estratégias.
  \item \textbf{Apoio adicional}: aluno com regressão ou estagnação recebe
        plano de apoio individualizado e, se necessário, encaminhamento
        ao PAEE.
\end{itemize}

\subsection{Acomodações Neuroinclusivas na Avaliação}

Para estudantes com TEA, dislexia ou outras condições neurodivergentes, a
avaliação pode ser adaptada sem prejuízo do que se pretende medir:

\begin{itemize}
  \item \textbf{Tempo flexível}: sem cronômetro; pausas livres a qualquer momento.
  \item \textbf{Instrução oral}: leitura dos enunciados pelo professor ou
        por leitor de tela, sem alterar o conteúdo.
  \item \textbf{Resposta oral}: o aluno pode ditar as respostas quando a
        escrita for uma barreira adicional.
  \item \textbf{Formato ampliado}: fonte maior e uma questão por página,
        quando necessário.
  \item \textbf{Redução de estímulos}: sala silenciosa, iluminação suave,
        sem elementos decorativos que desviem a atenção.
\end{itemize}

\checklistdominio{90\%+ de acertos com qualidade em avaliação diagnóstica}
